\documentclass{article}
\usepackage[utf8]{inputenc}
\usepackage[italian]{babel}
\usepackage{amsmath, amssymb}
\usepackage{tikz}

\begin{document}

\begin{flushright}
FISICA 1 - Fabrizio Cei \quad 27/02/2024
\end{flushright}

2) $\omega_1 \neq \omega_2$ ($\varphi_1 = \varphi_2 = 0$ per semplicità) (non armonico)
\begin{equation*}
x(t) = A_1 \cos(\omega_1 t) + A_2 \cos(\omega_2 t)
\end{equation*}
\begin{align*}
\ddot{x}(t) &= -\omega_1^2 A_1 \cos(\omega_1 t) - \omega_2^2 A_2 \cos(\omega_2 t) \\
&\neq -\lambda (A_1 \cos(\omega_1 t) + A_2 \cos(\omega_2 t)) = -\lambda x(t)
\end{align*}
$\nexists \lambda$ perché $\omega_1^2 \neq \omega_2^2$ quindi non si può scrivere una equazione ARMONICA.

\subsection*{2.1) $A_1 = A_2$ FENOMENO DEI BATTIMENTI (se $\omega_1 \simeq \omega_2$)}
\begin{align*}
x(t) &= A_1 [\cos(\omega_1 t) + \cos(\omega_2 t)] \\
&= 2A_1 \cdot \cos\left[ \left( \frac{\omega_1 + \omega_2}{2} \right) t \right] \cos\left[ \left( \frac{\omega_1 - \omega_2}{2} \right) t \right] \\
&= A(t) \cos\left[ \left( \frac{\omega_1 + \omega_2}{2} \right) t \right]
\end{align*}
con $A(t) = 2A_1 \cos\left[ \left( \frac{\omega_1 - \omega_2}{2} \right) t \right]$, chiamando $\Omega_2 = \frac{\omega_1 - \omega_2}{2}$.\\
Questa è un'OSCILLAZIONE ARMONICA con $\Omega_1 = (\omega_1 + \omega_2)/2$ e ampiezza $A(t)$ variabile.

\vspace{0.5cm}
\begin{minipage}{0.4\textwidth}
\begin{tikzpicture}[scale=0.8]
\draw[->] (-0.2,0) -- (5,0) node[right] {$t$};
\draw[->] (0,-2.5) -- (0,2.5) node[above] {$x$};
\node[left] at (0,2) {$2A_1$};
\node[left] at (0,-2) {$-2A_1$};

% Inviluppo
\draw[blue, thick, domain=0:4.5, samples=100] plot (\x, {2*cos(\x*180/2)});
\draw[blue, thick, domain=0:4.5, samples=100] plot (\x, {-2*cos(\x*180/2)});

% Onda portante
\draw[black, domain=0:4.5, samples=200] plot (\x, {2*cos(\x*180/2)*cos(\x*180*2)});

\node[below, font=\scriptsize] at (2,-2.5) {OSC. RAPIDA / OSC. LENTA};
\end{tikzpicture}
\end{minipage}
\hfill
\begin{minipage}{0.55\textwidth}
Si vengono a creare dei LOBI, generati dalla pulsazione armonica $\Omega_2$, detti ONDE MODULANTI che descrivono come varia nel tempo l'ampiezza dell'oscillazione con pulsazione $\Omega_1$ le cui curve sono dette ONDE PORTANTI e sono modulate da $\Omega_2$.
\end{minipage}

\vspace{0.5cm}
\begin{equation*}
T_2 = \frac{2\pi}{\Omega_2} = \frac{4\pi}{\omega_1 - \omega_2} \qquad T_1 = \frac{2\pi}{\Omega_1} = \frac{4\pi}{\omega_1 + \omega_2}
\end{equation*}

All'interno delle oscillazioni modulanti (in blu) i massimi (e i minimi) vengono raggiunti 2 volte (e 2 volte) mentre normalmente, all'interno di un periodo $T_2$ essi dovrebbero essere raggiunti 1 volta (e 1 volta).\\
Quindi la distanza temporale $\Delta t$ fra un massimo (e un minimo) e l'altro è pari a (BATTIMENTO):
\begin{equation*}
|\Delta t| = T_2 / 2 = \frac{2\pi}{|\omega_1 - \omega_2|} = \frac{1}{|\omega_1 / 2\pi - \omega_2 / 2\pi|} = \frac{1}{|f_1 - f_2|} = \frac{1}{|\Delta f|}
\end{equation*}
dove $\Delta f$ è la differenza fra la frequenza di $\omega_1$ e di $\omega_2$.

$\implies$ questo fenomeno avviene per esempio in musica quando si accordano gli strumenti col DIAPASON: se l'accordatore percepisce la modulazione di ampiezza dell'onda risultante dalla composizione del suono del diapason e della corrispondente corda dello strumento significa che i due suoni non sono identici, quindi agisce sul tasto dello strumento da accordare in modo da riportare in fase le due oscillazioni ($\omega_1 = \omega_2$).

\end{document}
